% S T R E S Z C Z E N I E
% ---------------

\section*{Streszczenie}
Praca dokumentuje powstanie gry \textit{Echoes of Vanitas}, czyli 2.5D RPG zbudowanego od podstaw w silniku Godot 4 i języku C\#. Akcja toczy się w świecie zamieszkałym przez antropomorficzne zwierzęta; głównym bohaterem jest Nietoperz Złodziej, który trafia na księgę spisaną w języku dawno wymarłej rasy ludzi. Tytuł powstał jako jednoosobowa produkcja indie: kod, sceny i dokumentacja techniczna są moją pracą.

Jednym z filarów projektu jest autorska warstwa graficzna. Całą oprawę, od szkiców koncepcyjnych po finalne sprite'y, postacie niegrywalne, ikony ekwipunku i elementy interfejsu, narysowałem własnoręcznie w programie Clip Studio Paint, bez gotowych paczek i generatorów. Spójny, ręcznie rysowany styl 2D łączy się tu z trójwymiarowym oświetleniem renderera \textit{Forward Plus}, dając charakterystyczną oprawę 2.5D.

Część implementacyjna obejmuje hub świata powierzchni z postaciami niegrywalnymi, proceduralnie generowany loch (trzy warianty układu różniące się liczbą pokoi i paletą świateł), losowane modyfikatory runu (więcej monet, dodatkowa elita w pokoju bossa, zaciemnienie), skrzynie z nagrodami, cztery rodzaje przeciwników (zwykli wrogowie, elity ze sztuczną inteligencją uczącą się reakcji gracza, boss z dashem i bombami obszarowymi oraz niewidzialni ujawniani dopiero przez radar), pięć mini-gier w hubie (układanie bloków, labirynt, sekwencja świateł, łamanie zamka, łączenie rur), system ekwipunku z ulepszaniem przedmiotów, sekwencyjny quest główny \textit{Pętla Vanitas} (klucz z mini-gry, loch, relikt, rozmowa ze strażnikiem), zapis stanu gry, minimapę z maską kołową renderowaną w \texttt{SubViewport} oraz pełny interfejs (HUD, menu zakładkowe, toasty, ekran startowy).

Kluczowym wkładem własnym jest \texttt{EliteBrain}, moduł adaptacyjnej AI dla elit oparty na kontekstowym bandycie wieloramiennym. Po zakończeniu starcia mózg ocenia, jak dobrze poradziła sobie wybrana taktyka (oskrzydlanie, szarża, utrzymywanie dystansu lub zasadzka), i koryguje wartości akcji $Q$ aktualizacją EMA. Selekcja taktyki używa polityki $\varepsilon$-zachłannej, a kontekstem są dwie zmienne dyskretyzowane: stan zdrowia gracza i liczba żywych elit. Równolegle prowadzony model gracza (uśredniona prędkość, tendencja do wycofywania się, DPS) wpływa na czas predykcji ruchu, dystans preferowany przez elity oraz parametry agresji i ostrożności. Stan mózgu zapisywany jest do wersjonowanego pliku JSON, dzięki czemu nauka utrzymuje się pomiędzy sesjami.

Pod względem architektury projekt opiera się na klasycznych wzorcach: maszynie stanów dla gracza i wrogów, singletonach dla usług globalnych (sześć autoloadów oraz \texttt{DungeonManager} i \texttt{UIController} w scenie głównej), obserwatorze w postaci statycznej klasy \texttt{GameEvents} oraz natywnym dla Godot wzorcu komponentowym. Loch generowany jest wariantem błądzenia losowego na siatce $9 \times 9$, z przypisaniem pokoi specjalnych (boss, skarb), spawnem wrogów o wspólnym balansie HP/siły i wypiekiem siatki nawigacyjnej. Radar w lochu używa półkulistego \texttt{Area3D} i odpytuje nakładające się ciała co klatkę fizyki; po trzech sekundach przebywania w polu wykrywania niewidzialny przeciwnik przełącza shader (\texttt{sprite\_visibility}) na pełną widoczność.

Wynik jest grywalny, kompletny technicznie i daje solidną podstawę pod dalszy rozwój, w szczególności rozbudowę narracji, dodanie kolejnych lokacji oraz eksperymenty z głębokim uczeniem wzmocnionym, do którego dyskretna tablica $Q$ jest tu jedynie pierwszym krokiem.

\cleardoublepage
